\documentclass[11pt]{article}

\usepackage[utf8]{inputenc}
\usepackage[T1]{fontenc}
\usepackage{lmodern}
\usepackage{geometry}
\usepackage{amsmath,amssymb,amsfonts,amsthm,mathtools}
\usepackage{bm}
\usepackage{enumitem}
\usepackage{microtype}
\usepackage{hyperref}
\usepackage{titlesec}
\usepackage{booktabs}
\usepackage{array}
\usepackage{setspace}
\usepackage{mathrsfs}
\usepackage{physics}

\geometry{margin=1in}
\hypersetup{colorlinks=true, linkcolor=black, urlcolor=blue, citecolor=black}
\setlist[enumerate]{topsep=0.4em,itemsep=0.35em}
\setlist[itemize]{topsep=0.3em,itemsep=0.25em}
\titleformat{\section}{\large\bfseries}{\thesection.}{0.5em}{}
\titleformat{\subsection}{\normalsize\bfseries}{\thesubsection.}{0.5em}{}
\setstretch{1.06}

\newcommand{\Aether}{\AE{}ther}
\newcommand{\Aflow}{\AE{}ther-flow}
\newcommand{\TheoryName}{\Aflow{} Interpretation of Relativity}
\newcommand{\TheoryProgram}{\Aether{} / \Aflow{} framework}

\newcommand{\CompletedMetric}{g^{\sharp}}
\newcommand{\MatterSpace}{\Psi}

\newcommand{\Car}{\mathrm{car}}
\newcommand{\Cur}{\mathrm{cur}}
\newcommand{\Hom}{\mathrm{hom}}

\newcommand{\ForSpace}[1]{\mathcal I_{#1}^{\mathrm{for}}}
\newcommand{\PrimShell}{\mathcal S_{\mathrm{prim}}}
\newcommand{\Reduction}[1]{\mathcal R_{#1}}
\newcommand{\CurrentCarrier}[1]{\mathfrak C_{#1}^{\mathrm{cur}}}
\newcommand{\EqCarrier}[1]{\mathfrak C_{#1}^{\mathrm{eq}}}
\newcommand{\MetricOut}{\mathscr G_{\mathrm{out}}}
\newcommand{\MatterOut}{\mathscr T_{\mathrm{out}}}

\newcommand{\SubTarget}[1]{E_{#1}^{\mathrm{sub}}}
\newcommand{\SubPair}[1]{\mathfrak s_{#1}^{\mathrm{sub}}}
\newcommand{\EqnTarget}[1]{Q_{#1}^{\mathrm{eqn}}}
\newcommand{\HiddenSpace}[1]{\mathcal H_{#1}^{\mathrm{hid}}}
\newcommand{\HiddenBody}[1]{D_{#1}^{\mathrm{hid}}}

\newcommand{\ProtoSpace}[1]{\mathcal U_{#1}^{\mathrm{proto}}}
\newcommand{\ProtoBody}[1]{\mathcal B_{#1}^{\mathrm{proto}}}
\newcommand{\ProtoSection}[1]{\Gamma_{#1}^{\mathrm{proto}}}
\newcommand{\ProtoShellSet}[1]{\mathcal Z_{#1}^{\mathrm{proto}}}
\newcommand{\ProtoZeroCore}[1]{\mathcal Z_{#1}^{0,\mathrm{proto}}}
\newcommand{\ProtoSplit}[1]{\Phi_{#1}^{\mathrm{proto}}}
\newcommand{\ProtoCharge}[1]{\mathcal Q_{#1}^{\mathrm{proto}}}
\newcommand{\ProtoEmbed}[1]{\zeta_{#1}^{\mathrm{proto}}}
\newcommand{\ProtoKernelCh}[1]{\widetilde K_{#1}^{0,\mathrm{proto}}}
\newcommand{\ProtoStateComp}[1]{P_{#1}^{\mathrm{proto,co}}}

\newcommand{\PreProtoNucleus}{\mathfrak Q^{\mathrm{nuc}}}
\newcommand{\PreProtoSpace}[1]{\mathcal V_{#1}^{\mathrm{preproto}}}
\newcommand{\PreProtoBody}[1]{\mathcal B_{#1}^{\mathrm{preproto}}}
\newcommand{\PreProtoProj}[1]{\Pi_{#1}^{\mathrm{preproto}}}
\newcommand{\PreProtoProtoMin}[1]{\mathfrak f_{#1}^{\mathrm{preproto,proto}}}
\newcommand{\PreProtoSection}[1]{\Gamma_{#1}^{\mathrm{preproto}}}
\newcommand{\PreProtoFunctional}[1]{\mathscr W_{#1}^{\mathrm{preproto}}}
\newcommand{\PreProtoShellSet}[1]{\mathcal Z_{#1}^{\mathrm{preproto}}}

\newcommand{\PreNucPackage}{\mathfrak R^{\mathrm{prenuc}}}
\newcommand{\PreNucSpace}[1]{\mathcal W_{#1}^{\mathrm{prenuc}}}
\newcommand{\PreNucBody}[1]{\mathcal B_{#1}^{\mathrm{prenuc}}}
\newcommand{\PreNucProj}[1]{\Pi_{#1}^{\mathrm{prenuc}}}
\newcommand{\PreNucNucleusMin}[1]{\mathfrak f_{#1}^{\mathrm{prenuc,nuc}}}
\newcommand{\PreNucSection}[1]{\Gamma_{#1}^{\mathrm{prenuc}}}
\newcommand{\PreNucFunctional}[1]{\mathscr W_{#1}^{\mathrm{prenuc}}}
\newcommand{\PreNucShellSet}[1]{\mathcal Z_{#1}^{\mathrm{prenuc}}}

\newtheorem{definition}{Definition}
\newtheorem{proposition}{Proposition}
\newtheorem{remark}{Remark}
\newtheorem{corollary}{Corollary}
\newtheorem{theorem}{Theorem}

\title{\textbf{The \TheoryName{}}\\[0.4em]
\large Primitive-Reservoir Observer-Reduction\\
\large Proto-Anchored Exact-Shell Transport Nucleus\\
\large from Nucleus-Generating Substrate Pre-Nucleus Dynamics}
\author{}
\date{}

\begin{document}

\maketitle

\begin{abstract}
The current primitive-reservoir observer-reduction frontier is the proto-anchored exact-shell transport nucleus. The previous collapse theorem already spent the remaining internal compression move available at that scope: the larger observer-relevant pre-proto-germ core carries no extra benchmark-relevant freedom beyond that smaller nucleus. After that result, the next honest benchmark-facing task is no longer another same-output pre-proto-germ relay or another restatement of the former observer-relevant core. It is either to derive or justify the nucleus from more primitive substrate structure, or to prove a genuinely smaller theorem-level collapse strictly inside it.

The present note takes the first route. For each sector it introduces \emph{nucleus-generating substrate pre-nucleus dynamics}: a compact convex pre-nucleus body over the recorded transport-nucleus state space, a unique pre-nucleus minimizer on the fixed transport-nucleus body, a smooth pre-nucleus restriction section, and an exact pre-nucleus shell projecting diffeomorphically to the exact transport-nucleus shell. The current proto-anchored exact-shell transport nucleus is then recovered canonically as the projected exact-shell transport image of that deeper pre-nucleus dynamics.

The benchmark boundary remains fixed throughout. The one operative metric is still exactly \(\CompletedMetric\), the ordinary matter route is unchanged, there is no second operative metric, the low-energy matter law is not enlarged, and no stronger promotion language is licensed. The positive result is therefore narrow but benchmark facing: the proto-anchored exact-shell transport nucleus is no longer left as an unexplained primitive stopping point on the active line. The remaining honest burden moves one layer deeper again, to the pre-nucleus dynamics themselves or to a sharper collapse inside the pre-nucleus package.
\end{abstract}

\section{Introduction}

The active derivational program of \TheoryName{} has now reached a cleaner burden than another same-output transport-nucleus restatement. The current transport-nucleus collapse theorem already proved that the benchmark-facing pre-proto-germ data collapse to a smaller proto-anchored exact-shell transport nucleus. The accompanying routing consequence is immediate: the next honest benchmark-facing manuscript must either derive or justify that nucleus from more primitive substrate structure, or prove a genuinely smaller collapse strictly inside it. Reopening the former larger observer-relevant core, re-running another same-output pre-proto-germ relay, or returning to stationary-reduction, stationary-Hessian, spectral, or selector layers would no longer address the live burden.

The present note takes the first route. Its positive claim is not that final substrate microphysics has been identified, and it is not that the benchmark exact-relativistic package has been changed. Its claim is narrower and benchmark facing. The current transport nucleus is shown to arise canonically from a more primitive pre-nucleus-level substrate dynamics whose projected exact-shell transport already contains the observer-relevant structure of the nucleus. In that sense the transport nucleus is no longer the primitive place where the active line has to stop.

Two features matter here. First, the theorem targets the current nucleus itself rather than revisiting the already-spent larger pre-proto-germ core, so it does not spend the next move on an older same-output presentation. Second, the deeper pre-nucleus package is not merely another renaming of the current nucleus. It adds an explicit pre-nucleus state space, pre-nucleus body, pre-nucleus-to-nucleus projection, and deeper exact-shell transport data below the current frontier. The current nucleus body, the proto anchor, the restriction section, and the exact shell are therefore projected exact-shell images of a deeper substrate package rather than unexplained primitive data.

\paragraph{Framework and claim boundary.}
\TheoryProgram{} denotes the broader \Aether{} / \Aflow{} framework built on the claims that the \Aether{} is the underlying four-dimensional substrate of reality, the \Aflow{} is its intrinsic ordered motion, observed three-dimensional space is the local experiential slice of that deeper substrate, S-time is the experienced order of change arising from matter, light, and the \Aflow{}, and observed expansion is the three-dimensional appearance of deeper four-dimensional ordered motion. Within that framework, \TheoryName{} denotes the exact-closure theory in which the effective gravitational dynamics are adopted to be exactly Einsteinian with universal matter coupling. In the active sequence, \emph{adoption} means use of that established relativistic dynamics without claiming substrate derivation, while \emph{derivation} is reserved for a first-principles recovery from explicit substrate variables.

The present note stays strictly inside that boundary. It does not alter the benchmark exact-closure package. It does not introduce a second operative metric or a new low-energy matter law. It only proves that the current proto-anchored exact-shell transport nucleus can itself be generated from a more primitive pre-nucleus-level substrate dynamics.

\section{Fixed Inputs}

\begin{definition}[Recorded primitive continuation data]
\label{def:fixed}
Fix the following continuation data from the active primitive-reservoir observer-reduction line.
\begin{enumerate}
    \item For each sector label \(\sigma\in\{\Car,\Cur,\Hom\}\), a primitive forcing shell
    \[
    \ForSpace{\sigma},
    \]
    a visible reduction map
    \[
    \Reduction{\sigma}:\ForSpace{\sigma}\to X_{\sigma},
    \]
    and shell-level current and equilibrium carriers
    \[
    \CurrentCarrier{\sigma}:\ForSpace{\sigma}\to Z_{\sigma}^{\mathrm{cur}},
    \qquad
    \EqCarrier{\sigma}:\ForSpace{\sigma}\to Z_{\sigma}^{\mathrm{eq}}.
    \]
    \item The product shell
    \[
    \PrimShell
    \equiv
    \ForSpace{\Car}\times \ForSpace{\Cur}\times \ForSpace{\Hom}.
    \]
    \item The recorded metric and ordinary-matter outputs
    \[
    \MetricOut:\PrimShell\to\{\text{observer metrics}\},
    \qquad
    \MatterOut:\PrimShell\times\MatterSpace\to\{\text{observer stress tensors}\},
    \]
    satisfying
    \begin{equation}
    \MetricOut(\mathbf w)=\CompletedMetric,
    \qquad
    \MatterOut(\mathbf w,\psi)
    =
    T_{\mu\nu}^{\mathrm{matter}}[\CompletedMetric,\psi]
    \label{eq:fixedoutputs}
    \end{equation}
    for every \(\mathbf w\in\PrimShell\) and every matter configuration \(\psi\in\MatterSpace\).
\end{enumerate}
\end{definition}

\begin{definition}[Recorded proto-germ exact-shell target]
\label{def:proto}
Fix \(\sigma\in\{\Car,\Cur,\Hom\}\). The active line already fixes:
\begin{enumerate}
    \item the same substrate target and pairing
    \[
    \SubTarget{\sigma},
    \qquad
    \SubPair{\sigma}:
    \SubTarget{\sigma}\times\SubTarget{\sigma}\to\mathbb R;
    \]
    \item a hidden fiber space \(\HiddenSpace{\sigma}\) with compact hidden body \(\HiddenBody{\sigma}\subset\HiddenSpace{\sigma}\);
    \item a finite-dimensional proto-germ state space \(\ProtoSpace{\sigma}\), a compact convex proto-germ body \(\ProtoBody{\sigma}\subset\ProtoSpace{\sigma}\), a smooth proto-germ restriction section
    \[
    \ProtoSection{\sigma}:\ProtoSpace{\sigma}\to\SubTarget{\sigma},
    \]
    and the exact proto-germ shell
    \[
    \ProtoShellSet{\sigma}
    \equiv
    \bigl(\ProtoSection{\sigma}\bigr)^{-1}(0)\cap\ProtoBody{\sigma};
    \]
    \item a closed embedded proto zero core \(\ProtoZeroCore{\sigma}\subset\ProtoShellSet{\sigma}\) and a split diffeomorphism
    \[
    \ProtoSplit{\sigma}:
    \ProtoZeroCore{\sigma}\times\HiddenBody{\sigma}
    \xrightarrow{\cong}
    \ProtoShellSet{\sigma};
    \]
    \item a hidden charge
    \[
    \ProtoCharge{\sigma}:\HiddenSpace{\sigma}\to\EqnTarget{\sigma},
    \]
    a smooth observer-compatible exact proto-germ zero-response realization
    \[
    \ProtoEmbed{\sigma}:\ForSpace{\sigma}\hookrightarrow\ProtoZeroCore{\sigma},
    \]
    and a fixed-data solvability / coercive package on \(\ProtoZeroCore{\sigma}\), recorded through \(\ProtoKernelCh{\sigma}\) and \(\ProtoStateComp{\sigma}\).
\end{enumerate}
\end{definition}

\begin{remark}
Definitions \ref{def:fixed} and \ref{def:proto} are fixed inputs. The one operative metric, the ordinary matter route, and the recorded proto-germ exact-shell target are not reopened here. The present note asks only whether the current proto-anchored exact-shell transport nucleus can itself be generated from more primitive substrate dynamics.
\end{remark}

\section{Recorded Current Transport-Nucleus Target}

\begin{definition}[Recorded proto-anchored exact-shell transport nucleus]
\label{def:nucleus}
Fix \(\sigma\in\{\Car,\Cur,\Hom\}\). The recorded current transport nucleus \(\PreProtoNucleus\) for sector \(\sigma\) consists of:
\begin{enumerate}
    \item the pre-proto-germ state space \(\PreProtoSpace{\sigma}\), compact convex body \(\PreProtoBody{\sigma}\subset\PreProtoSpace{\sigma}\), projection
    \[
    \PreProtoProj{\sigma}:\PreProtoSpace{\sigma}\to\ProtoSpace{\sigma},
    \]
    and proto section
    \[
    \PreProtoProtoMin{\sigma}:\ProtoBody{\sigma}\to\PreProtoBody{\sigma}
    \]
    satisfying
    \[
    \PreProtoProj{\sigma}\bigl(\PreProtoBody{\sigma}\bigr)=\ProtoBody{\sigma},
    \qquad
    \PreProtoProj{\sigma}\circ\PreProtoProtoMin{\sigma}
    =
    \mathrm{id}_{\ProtoBody{\sigma}};
    \]
    \item the restriction section
    \[
    \PreProtoSection{\sigma}:\PreProtoSpace{\sigma}\to\SubTarget{\sigma},
    \]
    the induced functional
    \[
    \PreProtoFunctional{\sigma}(q)
    \equiv
    \SubPair{\sigma}\bigl(\PreProtoSection{\sigma}(q),\PreProtoSection{\sigma}(q)\bigr),
    \]
    and the exact shell
    \[
    \PreProtoShellSet{\sigma}
    \equiv
    \bigl(\PreProtoSection{\sigma}\bigr)^{-1}(0)\cap\PreProtoBody{\sigma},
    \]
    with
    \[
    \PreProtoSection{\sigma}\circ\PreProtoProtoMin{\sigma}
    =
    \ProtoSection{\sigma},
    \qquad
    \PreProtoProj{\sigma}\big|_{\PreProtoShellSet{\sigma}}
    :
    \PreProtoShellSet{\sigma}
    \xrightarrow{\cong}
    \ProtoShellSet{\sigma};
    \]
    \item benchmark faithfulness, meaning preservation of the benchmark outputs \eqref{eq:fixedoutputs}, no second operative metric, no enlarged low-energy matter law, and no premature promotion from adoption to completed derivation.
\end{enumerate}
By the already-recorded current frontier theorem, the zero core, split, hidden charge, exact zero-response realization, and fixed-data solvability / coercive package are then canonically induced downstream from this nucleus and the fixed proto-germ target rather than taken as additional independent benchmark-facing data.
\end{definition}

\begin{remark}
Definition \ref{def:nucleus} is the recorded target already fixed by the current frontier theorem. The present note does not spend its move on a second restatement of that target. It records it only because the deeper pre-nucleus dynamics are defined relative to it.
\end{remark}

\begin{remark}[Downstream fixed fact]
\label{rem:downstream}
The already-recorded current frontier theorem proves that once a proto-anchored exact-shell transport nucleus is fixed, the larger observer-relevant pre-proto-germ core is canonically recovered and carries no extra benchmark-relevant freedom. The present note uses that downstream fact only as a routing consequence: after the induced nucleus is obtained, no second larger current-core package needs to be postulated anew.
\end{remark}

\section{Nucleus-Generating Substrate Pre-Nucleus Dynamics}

\begin{definition}[Nucleus-generating substrate pre-nucleus dynamics]
\label{def:prenuc}
Fix \(\sigma\in\{\Car,\Cur,\Hom\}\). A nucleus-generating substrate pre-nucleus dynamics package \(\PreNucPackage\) for sector \(\sigma\) consists of the following data.
\begin{enumerate}
    \item A finite-dimensional Euclidean pre-nucleus state space \(\PreNucSpace{\sigma}\), a compact convex pre-nucleus body
    \[
    \PreNucBody{\sigma}\subset\PreNucSpace{\sigma}
    \]
    with nonempty interior, an affine surjection
    \[
    \PreNucProj{\sigma}:\PreNucSpace{\sigma}\to\PreProtoSpace{\sigma},
    \]
    and a smooth pre-nucleus fiber minimizer
    \[
    \PreNucNucleusMin{\sigma}:\PreProtoBody{\sigma}\to\PreNucBody{\sigma}
    \]
    satisfying
    \[
    \PreNucProj{\sigma}\bigl(\PreNucBody{\sigma}\bigr)=\PreProtoBody{\sigma},
    \qquad
    \PreNucProj{\sigma}\circ\PreNucNucleusMin{\sigma}
    =
    \mathrm{id}_{\PreProtoBody{\sigma}}.
    \]
    \item A smooth pre-nucleus restriction section
    \[
    \PreNucSection{\sigma}:\PreNucSpace{\sigma}\to\SubTarget{\sigma},
    \]
    the pre-nucleus functional
    \[
    \PreNucFunctional{\sigma}(w)
    \equiv
    \SubPair{\sigma}\bigl(\PreNucSection{\sigma}(w),\PreNucSection{\sigma}(w)\bigr),
    \]
    and the exact pre-nucleus shell
    \[
    \PreNucShellSet{\sigma}
    \equiv
    \bigl(\PreNucSection{\sigma}\bigr)^{-1}(0)\cap\PreNucBody{\sigma},
    \]
    satisfying
    \[
    \PreNucSection{\sigma}\circ\PreNucNucleusMin{\sigma}
    =
    \PreProtoSection{\sigma},
    \qquad
    \PreNucProj{\sigma}\big|_{\PreNucShellSet{\sigma}}
    :
    \PreNucShellSet{\sigma}
    \xrightarrow{\cong}
    \PreProtoShellSet{\sigma}.
    \]
\end{enumerate}
The package is \emph{benchmark faithful} if it preserves the benchmark outputs \eqref{eq:fixedoutputs}, introduces no second operative metric, introduces no enlarged low-energy matter law, and is presented as a deeper substrate-side source for the current proto-anchored exact-shell transport nucleus rather than as a replacement benchmark theory.
\end{definition}

\begin{remark}
Definition \ref{def:prenuc} is intentionally narrower than another full pre-proto-germ relay. It does not rebuild the former observer-relevant pre-proto-germ core as an independent datum. It only introduces the deeper state space, body, and exact-shell transport data needed to generate the current nucleus.
\end{remark}

\begin{remark}
The label \emph{pre-nucleus} is used here only as a local marker for a deeper layer below the current transport-nucleus frontier. It does not by itself assert a completed microscopic ontology or a final primitive law. Its role is only to name the next sharper transport package below the current nucleus.
\end{remark}

\section{Projected Exact-Shell Transport to the Current Nucleus}

\begin{proposition}[The proto anchor persists one layer deeper]
\label{prop:anchor}
Fix a benchmark-faithful pre-nucleus package. Then
\[
(\PreProtoProj{\sigma}\circ\PreNucProj{\sigma})\bigl(\PreNucBody{\sigma}\bigr)
=
\ProtoBody{\sigma},
\]
\[
(\PreProtoProj{\sigma}\circ\PreNucProj{\sigma})\circ
\PreNucNucleusMin{\sigma}\circ\PreProtoProtoMin{\sigma}
=
\mathrm{id}_{\ProtoBody{\sigma}},
\]
\[
\PreNucSection{\sigma}\circ\PreNucNucleusMin{\sigma}\circ\PreProtoProtoMin{\sigma}
=
\ProtoSection{\sigma},
\]
and the composite shell transport
\[
(\PreProtoProj{\sigma}\circ\PreNucProj{\sigma})\big|_{\PreNucShellSet{\sigma}}
:
\PreNucShellSet{\sigma}
\xrightarrow{\cong}
\ProtoShellSet{\sigma}
\]
is a diffeomorphism.
\end{proposition}

\begin{proof}
Definition \ref{def:prenuc}(1) gives
\[
\PreNucProj{\sigma}\bigl(\PreNucBody{\sigma}\bigr)=\PreProtoBody{\sigma}
\]
and
\[
\PreNucProj{\sigma}\circ\PreNucNucleusMin{\sigma}
=
\mathrm{id}_{\PreProtoBody{\sigma}}.
\]
Applying Definition \ref{def:nucleus}(1) to those identities yields the first two displayed formulas. Likewise,
\[
\PreNucSection{\sigma}\circ\PreNucNucleusMin{\sigma}
=
\PreProtoSection{\sigma}
\]
by Definition \ref{def:prenuc}(2), and
\[
\PreProtoSection{\sigma}\circ\PreProtoProtoMin{\sigma}
=
\ProtoSection{\sigma}
\]
by Definition \ref{def:nucleus}(2), so the third displayed formula follows by composition.

For the shell transport, Definition \ref{def:prenuc}(2) gives a diffeomorphism
\[
\PreNucProj{\sigma}\big|_{\PreNucShellSet{\sigma}}
:
\PreNucShellSet{\sigma}
\xrightarrow{\cong}
\PreProtoShellSet{\sigma},
\]
while Definition \ref{def:nucleus}(2) gives a diffeomorphism
\[
\PreProtoProj{\sigma}\big|_{\PreProtoShellSet{\sigma}}
:
\PreProtoShellSet{\sigma}
\xrightarrow{\cong}
\ProtoShellSet{\sigma}.
\]
Their composition is therefore a diffeomorphism from \(\PreNucShellSet{\sigma}\) to \(\ProtoShellSet{\sigma}\).
\end{proof}

\begin{proposition}[Pre-nucleus dynamics canonically reproduce the current transport nucleus]
\label{prop:nucleus}
Fix a benchmark-faithful pre-nucleus package. Then projected exact-shell transport along
\[
\PreNucProj{\sigma}:\PreNucSpace{\sigma}\to\PreProtoSpace{\sigma}
\]
reproduces exactly the recorded proto-anchored exact-shell transport nucleus of Definition \ref{def:nucleus}.
\end{proposition}

\begin{proof}
The current nucleus body and its nucleus-level sectioning data are recovered by Definition \ref{def:prenuc}(1): the pre-nucleus body projects onto the recorded current body, and the pre-nucleus nucleus minimizer sections that projection on \(\PreProtoBody{\sigma}\). The recorded proto anchor then persists by Proposition \ref{prop:anchor}.

The current nucleus restriction data are recovered by Definition \ref{def:prenuc}(2): the recorded current section is the minimizing transport of the pre-nucleus section, the current functional is therefore the projected restriction functional, and the exact current shell is the projected image of the exact pre-nucleus shell under a diffeomorphism.

Those are precisely the constituents retained by the current transport-nucleus frontier. Therefore the recorded current nucleus is not postulated independently inside this deeper class. It is the projected exact-shell transport image of the pre-nucleus package.
\end{proof}

\begin{proposition}[The benchmark package remains fixed]
\label{prop:boundary}
For every benchmark-faithful pre-nucleus package, the operative metric remains exactly \(\CompletedMetric\), the ordinary matter route remains unchanged, no second operative metric is introduced, and the low-energy matter law is not enlarged.
\end{proposition}

\begin{proof}
This is part of the benchmark-faithfulness requirement in Definition \ref{def:prenuc}. Equation \eqref{eq:fixedoutputs} remains unchanged on the full primitive forcing shell, so the induced nucleus does not alter the benchmark exact-closure package.
\end{proof}

\begin{proposition}[The larger current core is still recovered downstream]
\label{prop:downstream}
Fix a benchmark-faithful pre-nucleus package. Once projected exact-shell transport reproduces the current nucleus, the already-recorded transport-nucleus theorem canonically recovers the full observer-relevant pre-proto-germ core without any new larger datum.
\end{proposition}

\begin{proof}
Proposition \ref{prop:nucleus} recovers the recorded current nucleus. By Remark \ref{rem:downstream}, the already-recorded current frontier theorem then applies verbatim and canonically supplies the zero core, split, hidden charge, exact zero-response realization, and fixed-data solvability / coercive package from that nucleus together with the fixed proto-germ target. No second independently chosen larger current-core datum is added here.
\end{proof}

\section{Main Theorem}

\begin{theorem}[Proto-anchored exact-shell transport nucleus from nucleus-generating substrate pre-nucleus dynamics]
\label{thm:main}
Fix the primitive continuation data of Definition \ref{def:fixed}, the recorded proto-germ exact-shell target of Definition \ref{def:proto}, and the recorded current transport nucleus of Definition \ref{def:nucleus}. Then, for each sector \(\sigma\in\{\Car,\Cur,\Hom\}\), every benchmark-faithful nucleus-generating substrate pre-nucleus dynamics package canonically reproduces the current proto-anchored exact-shell transport nucleus with the same benchmark one-metric outputs. More precisely:
\begin{enumerate}
    \item projected exact-shell transport along \(\PreNucProj{\sigma}\) reproduces exactly the recorded current nucleus by Proposition \ref{prop:nucleus};
    \item the current transport nucleus is therefore no longer an unexplained primitive stopping point on the active line, but the projected exact-shell transport image of a more primitive pre-nucleus dynamics;
    \item the benchmark boundary remains fixed by Proposition \ref{prop:boundary};
    \item by Proposition \ref{prop:downstream}, the full observer-relevant pre-proto-germ core is then recovered downstream from that induced nucleus without any new larger datum;
    \item consequently the remaining honest burden moves one layer deeper again: derive or justify the nucleus-generating substrate pre-nucleus dynamics themselves from still more primitive structure, or else prove a still sharper collapse inside the pre-nucleus package.
\end{enumerate}
\end{theorem}

\begin{proof}
Item (1) is Proposition \ref{prop:nucleus}. Item (2) is the project-level consequence of item (1): once the pre-nucleus package is fixed, the current transport nucleus is no longer introduced as an independent primitive tuple. It is the projected exact-shell transport image of a deeper state space, deeper body, deeper section, and deeper exact shell.

Item (3) is Proposition \ref{prop:boundary}. Item (4) is Proposition \ref{prop:downstream}. Item (5) is the routing consequence of the previous items together with the active safeguard against counting same-output deeper relays as new front-line progress once the live observer-relevant datum is unchanged.
\end{proof}

\begin{corollary}[What must be done next]
\label{cor:next}
After Theorem \ref{thm:main}, the next honest benchmark-facing manuscript must either:
\begin{enumerate}
    \item derive or justify the nucleus-generating substrate pre-nucleus dynamics from still more primitive substrate structure in a way that changes benchmark-facing status;
    \item identify a genuinely smaller theorem-level observer-relevant core strictly inside the pre-nucleus package and prove a sharper collapse to it;
    \item do either of the previous two while preserving the same benchmark one-metric package and the same claim boundary.
\end{enumerate}
Another same-output restatement of the current proto-anchored exact-shell transport nucleus, another same-output deeper-origin relay that leaves that nucleus unchanged, another reopening of the former observer-relevant pre-proto-germ core as if it were still the live burden, a reopening of older screened shell, lift, support, reservoir, ground, stationary-reduction, stationary-Hessian, spectral, or selector layers, or any move that introduces a second operative metric, enlarges the ordinary matter law, or uses stronger promotion language does not satisfy the present burden.
\end{corollary}

\begin{proof}
Theorem \ref{thm:main} removes the current transport nucleus as the unexplained stopping point on the active line. Therefore another manuscript that merely restates that same nucleus, or merely returns to an already-screened larger relay, does not advance the benchmark-facing status of the project. What remains open is either to go one honest layer deeper again, to the pre-nucleus dynamics themselves, or to discover a smaller theorem-level observer-relevant core inside the pre-nucleus package.
\end{proof}

\section{Conclusion}

The positive result of this note is narrow but benchmark facing. The current transport-nucleus collapse theorem had already shown that the larger observer-relevant pre-proto-germ core contains no extra benchmark-relevant freedom beyond a smaller proto-anchored exact-shell transport nucleus. The present theorem then removes the next unexplained stopping point by showing that this nucleus itself is the projected exact-shell transport image of a more primitive nucleus-generating substrate pre-nucleus dynamics package.

The scope remains disciplined. The benchmark package is unchanged: the one operative metric is still exactly \(\CompletedMetric\), the ordinary matter route is unchanged, there is no second operative metric, and the low-energy matter law is not enlarged. Nothing here upgrades adoption to completed derivation or licenses stronger promotion language.

The open burden therefore moves one honest step further again. The next benchmark-facing manuscript should derive or justify the pre-nucleus dynamics from still more primitive substrate structure, unless the sharper honest gain available instead is a still smaller theorem-level collapse inside the pre-nucleus package. Until then, the current result should be read for what it is: a deeper transport-nucleus derivation on the active line of \TheoryName{}, not a claim that the full first-principles program is finished.

\end{document}
